\documentclass{article}
\usepackage{arxiv}
%\usepackage{amsthm}
\usepackage{amssymb}
\usepackage[utf8]{inputenc} % allow utf-8 input
\usepackage[T1]{fontenc}    % use 8-bit T1 fonts
\usepackage{mathtools}
\usepackage{hyperref}       % hyperlinks
\usepackage{url}            % simple URL typesetting
\usepackage{booktabs}       % professional-quality tables
\usepackage{amsfonts}       % blackboard math symbols
\usepackage{nicefrac}       % compact symbols for 1/2, etc.
\usepackage{microtype}      % microtypography
\usepackage{graphicx}
\usepackage{natbib}
\usepackage{doi}
\usepackage{amsmath}
\usepackage{physics}
\usepackage{wrapfig}
\usepackage{bm}
\usepackage{algorithm,algpseudocode}
\usepackage{xcolor}
\usepackage{forest}
\usepackage{tcolorbox}
\usepackage{tikz-cd}
\usepackage{enumitem}

\newenvironment{proof}{\paragraph{Proof:}}{\hfill$\square$}
\newenvironment{solution}{\paragraph{Solution:}}{\hfill$\square$}

%—– Blackboard‐bold sets —–
\newcommand{\R}{\mathbb{R}}
\newcommand{\N}{\mathbb{N}}
\newcommand{\Z}{\mathbb{Z}}
\newcommand{\Q}{\mathbb{Q}}
\newcommand{\C}{\mathbb{C}}

%—– Probability & expectation —–
\newcommand{\Prb}{\mathbb{P}}
\newcommand{\E}{\mathbb{E}}
\newcommand{\Var}{\mathrm{Var}}
\newcommand{\Cov}{\mathrm{Cov}}

%—– Common math operators —–
\newcommand{\diag}{\operatorname{diag}}
\newcommand{\supp}{\operatorname{supp}}
\newcommand{\argmax}{\operatorname*{arg\,max}}
\newcommand{\argmin}{\operatorname*{arg\,min}}
\newcommand{\ind}{\mathbf{1}}

%—– Norms, inner products —–
\newcommand{\inner}[2]{\left\langle #1, #2 \right\rangle}

%—– Calligraphic sets —–
\newcommand{\cA}{\mathcal{A}}
\newcommand{\cS}{\mathcal{S}}
\newcommand{\cT}{\mathcal{T}}
\newcommand{\cF}{\mathcal{F}}
\newcommand{\cE}{\mathcal{E}}

%—– Miscellaneous —–
\newcommand{\eps}{\varepsilon}
\newcommand{\dt}{\mathrm{d}t}
\newcommand{\given}{\,\bigl\vert\,}

\newcommand{\half}{\tfrac12}
\newcommand{\thalf}{\tfrac32}

\newcommand{\todo}[1]{\textcolor{red}{[TODO: #1]}}

\newcommand{\note}[1]{\begin{tcolorbox}
\textbf{Note:} #1\end{tcolorbox}}

\newtheorem{theorem}{Theorem}
\newtheorem{proposition}[theorem]{Proposition}
\newtheorem{lemma}[theorem]{Lemma}
\newtheorem{corollary}[theorem]{Corollary}
\newtheorem{definition}[theorem]{Definition}
\newtheorem{assumption}[theorem]{Assumption}
\newtheorem{remark}[theorem]{Remark}
\newtheorem{example}[theorem]{Example}
\newtheorem{conjecture}[theorem]{Conjecture}

\title{Dyadic Multi-Channel Primal--Dual Reconstruction for Limited-Angle CT}
\date{}

\author{ \href{https://orcid.org/0000-0003-1700-756X}{William Chang} \\
    Department of Applied Mathematics\\
    University of California, Los Angeles\\
    Los Angeles, CA \\
    \url{https://williamc.me/} \\
    \And
}

\renewcommand{\headeright}{}
\renewcommand{\undertitle}{}

\hypersetup{
pdftitle={Dyadic Multi-Channel Primal-Dual CT Reconstruction},
pdfsubject={Limited-angle tomography},
pdfauthor={William Chang},
pdfkeywords={Chambolle-Pock, fan-beam CT, limited angle, multichannel preconditioning},
}

\begin{document}
\maketitle

\begin{abstract}
We propose a dyadic multi-channel extension of the Chambolle--Pock (CP)
primal--dual algorithm for limited-angle fan-beam CT reconstruction with
directional total variation (DTV) regularization. The ramp-filtered data
fidelity is split into a partition-of-unity bank of $k$ band-pass
channels whose supports halve in bandwidth toward DC, mirroring the
location of the limited-angle null space in Fourier space. Each band
carries an independent dual variable with its own step size $\sigma_i$
and data tolerance $\eps_i$, enabling preferential refinement of the
low-frequency content that is poorly determined by sparse-view data.
We derive the analogue of the standard CP convergence condition for the
$k$-channel product-space saddle point, give an analytic criterion for
choosing the dyadic depth $k$ from the detector geometry, and specify a
geometric $\sigma_i = 2^i \sigma_0$ scaling that compensates the
ramp's spectral bias so the operator norm grows only linearly in $k$
rather than exponentially.
\end{abstract}

%=====================================================================
\section{Introduction}
%=====================================================================

In sparse-view or limited-angle CT, the discrete X-ray transform $X$
has a sizable null space concentrated at low spatial frequencies, where
missing projection angles fail to constrain slowly-varying image
content \citep{sidky2012}. Single-channel iterative methods such as
Chambolle--Pock with TV regularization \citep{chambolle2011,sidky2012}
treat the data residual as a single object in the dual space and apply
one step size $\sigma$ to its entire spectral content. This is
suboptimal: the high-frequency portion of the residual converges
rapidly because it is well determined by the data, while the
low-frequency portion converges slowly precisely because it is the part
the prior must regularize. A single $\sigma$ cannot accommodate both
regimes.

The classical remedy is a two-channel split that separates the
data fidelity into low-frequency (LF) and high-frequency (HF) bands
\citep{sidky2014}. Each band gets its own dual variable and its own
step size, and the convergence condition becomes a product-space CP
inequality with diagonal dual preconditioning. In practice the two-channel
split provides only a modest improvement over single-channel CP, in
part because a binary split is too coarse: the LF tail that drives the
null-space ambiguity is concentrated in just the bottom few percent of
the spectrum, and a two-channel split that puts everything below some
cutoff $f_c$ into one LF channel still imposes a uniform $\sigma_\text{lo}$
across an unbalanced range.

We extend the two-channel idea to a \emph{dyadic} bank of $k$ channels,
where channel $0$ covers the top half of the spectrum, channel $i$
covers the dyadic shell $[f_\text{max}/2^{i+1},\, f_\text{max}/2^i]$,
and channel $k-1$ is the closed LF tail. The construction has three
desirable properties:
\begin{enumerate}[leftmargin=*]
  \item \textbf{Partition of unity.} By construction the squared
    filter magnitudes telescope to the single-channel ramp,
    $\sum_i |f_i(u)|^2 = |u| = |f_\text{single}(u)|^2$. Splitting does
    not inflate the operator norm relative to a partition-of-unity
    decomposition with matched $\sigma_i$.
  \item \textbf{Analytic depth selection.} The dyadic depth $k$ is
    determined by the detector geometry: the smallest band must
    contain at least $\sim 8$ Fourier bins for the Hanning shape to
    be discretely faithful, giving $k_\text{max} = \lfloor \log_2(\text{nbins}/8) \rfloor$.
  \item \textbf{LF-emphasized dual step schedule.} The choice
    $\sigma_i = 2^i \sigma_0$ scales the dual step inversely to the
    band's ramp energy, so $\sigma_i \|R_i X\|^2 \approx \sigma_0 \|R_0 X\|^2$
    is approximately constant across $i$. The resulting weighted
    operator norm grows linearly in $k$, $\|\Sigma^{1/2} K\|^2 \approx k \cdot \|R_0 X\|^2$,
    rather than exponentially.
\end{enumerate}

Each band's data residual is constrained to a closed $\ell_2$ ball of
radius $\eps_i \sqrt{\text{nrays}}$, giving $k$ independent feasibility
tolerances that can encode prior knowledge about per-band noise levels
or trust. The remainder of the variational problem---directional
TV in $x$ and $y$ via DTV duals, positivity, and an $\ell_1$ image
prior---is handled exactly as in the single-channel case
\citep{rigie2015,kim2017}.

%=====================================================================
\section{Preliminaries}
%=====================================================================

\subsection{Geometry and operators}

Let $f \in \R^{N}$ denote a discretized image on an $n_x \times n_y$
grid ($N = n_x n_y$). Let $X : \R^N \to \R^M$ be the fan-beam X-ray
transform discretized over $n_\text{views}$ source positions and
$n_\text{bins}$ detector bins ($M = n_\text{views} \cdot n_\text{bins}$).
The detector coordinate is $u \in [-L/2,\, L/2]$ with $L = \text{detectorlength}$
and $b_{00} = -L/2$. Let $\nabla_x, \nabla_y : \R^N \to \R^N$ be the
forward finite-difference gradient operators, with adjoints $\nabla_x^*$,
$\nabla_y^*$.

We work in the convention of \citet{sidky2014} in which a ramp-filter
$R = X \cdot F$ is incorporated into the data fidelity, where
$F : \R^M \to \R^M$ is the Fourier multiplier with profile
$\hat F(u) = \sqrt{|u|}$. The full filtered projector is $R X$, and we
write $R^* R X = X^* F^* F X$ as the filtered backprojection composed
with projection.

\subsection{Single-channel CP reconstruction}

The single-channel Chambolle--Pock saddle-point problem is
\begin{equation}\label{eq:cp-single}
  \min_{f \ge 0}\ \alpha\, \mathrm{TV}_\text{D}(f)
    + \beta\, \|f\|_1
    \quad \text{s.t.}\quad
    \|R(X f - g)\|_2 \le \eps \sqrt{M},
\end{equation}
where $g \in \R^M$ is the measured sinogram, $\mathrm{TV}_\text{D}$ is
the directional total variation \citep{rigie2015}, and the data
constraint is a closed $\ell_2$ ball of radius $\eps \sqrt{M}$ around
the data. The dual variables are $y \in \R^M$ for the data constraint,
$p_x, p_y \in \R^N$ for the gradient channels, and $q \in \R^N$ for
$\|f\|_1$. The standard CP iteration with relaxation $\rho$ produces
the primal--dual update with convergence condition
\begin{equation}\label{eq:cp-single-cond}
  \tau\, \sigma\, \|K\|_2^2 < 1,\qquad
  K = \begin{bmatrix} R X \\ \nabla_x \\ \nabla_y \\ I \end{bmatrix},
\end{equation}
where $\sigma$ is the dual step and $\tau$ the primal step
\citep{chambolle2011}. In practice $\tau\,\sigma\,\|K\|^2 = (1+\delta)^{-1}$
for a small slack $\delta > 0$.

\subsection{Two-channel split}

Let $H_c(u) = \tfrac12 (1 + \cos(\pi u/u_c))\, \ind_{|u|<u_c}$ be a
Hanning window with half-width $u_c = |b_{00}|/c$. The two-channel
split of \citet{sidky2014} writes the ramp filter as
\begin{equation}\label{eq:two-channel-filters}
  f_\text{hi}(u) = \sqrt{|u|}\,\sqrt{1 - H_{c_\text{hi}}(u)},
  \qquad
  f_\text{lo}(u) = \sqrt{|u|}\,\sqrt{H_{c_\text{lo}}(u)},
\end{equation}
and produces a saddle-point problem with independent duals
$y_\text{hi}, y_\text{lo}$. When $c_\text{hi} = c_\text{lo}$ this is a
partition of unity ($|f_\text{hi}|^2 + |f_\text{lo}|^2 = |u|$); when
$c_\text{hi} \neq c_\text{lo}$ the bands either overlap or leave a
coverage gap, both of which inflate $\|K\|$.

The convergence condition becomes the product-space CP inequality with
diagonal dual preconditioning $\Sigma = \diag(\sigma_\text{hi} I, \sigma_\text{lo} I)$:
\begin{equation}\label{eq:two-cond}
  \tau\, \left\|\Sigma^{1/2}
    \begin{bmatrix} R_\text{hi} X \\ R_\text{lo} X \end{bmatrix}
    \right\|_2^2 < 1,
\end{equation}
or equivalently $\tau\, \big(\sigma_\text{hi}\|R_\text{hi}X\|^2
+ \sigma_\text{lo}\|R_\text{lo}X\|^2\big) < 1$.

%=====================================================================
\section{Dyadic Multi-Channel Algorithm}
%=====================================================================

\subsection{Dyadic filter bank}\label{sec:dyadic-bank}

Fix $k \ge 2$. Define the dyadic Hanning sequence
$H_i(u) = H_{2^i}(u)$ for $i = 0, 1, \ldots, k-1$, with half-widths
$|b_{00}|/2^i$. The $k$ band-pass filters are defined by telescoping
differences:
\begin{equation}\label{eq:dyadic-filters}
  f_i(u) = \sqrt{|u|}\,\sqrt{w_i(u)},
  \qquad
  w_i(u) = \begin{cases}
    1 - H_1(u) & i = 0 \\
    H_i(u) - H_{i+1}(u) & 0 < i < k-1 \\
    H_{k-1}(u) & i = k-1.
  \end{cases}
\end{equation}
Channel $0$ is the top half of the spectrum, channel $i$
($0 < i < k-1$) is the dyadic shell $[|b_{00}|/2^{i+1},\,|b_{00}|/2^i]$,
and channel $k-1$ is the closed LF tail $[0,\,|b_{00}|/2^{k-1}]$. Each
$R_i = X \cdot F_i$ is the projector composed with the Fourier
multiplier of profile $f_i$. By construction $\sum_i w_i(u) = 1$ for
all $u$, so the squared filter magnitudes form a partition of unity:
\begin{equation}\label{eq:pou}
  \sum_{i=0}^{k-1} |f_i(u)|^2 = |u| = |f_\text{single}(u)|^2.
\end{equation}

\subsection{Analytic depth selection}\label{sec:depth}

The dyadic depth $k$ is bounded above by the detector's spectral
resolution: the lowest band has half-width $|b_{00}|/2^{k-1}$, which
corresponds to approximately $\text{nbins}/2^k$ Fourier bins on the
discrete grid. To keep the Hanning shape numerically faithful we
require at least $m \ge 8$ bins per band:
\begin{equation}\label{eq:k-max}
  k_\text{max} = \left\lfloor \log_2(\text{nbins}/m) \right\rfloor.
\end{equation}
For our $\text{nbins} = 1024$ detector with $m = 8$, this gives
$k_\text{max} = 7$. Deeper splits put the LF tail inside the
discretization noise floor and provide no further refinement of the
low-frequency dual variable.

\begin{remark}
The image-grid criterion---the lowest band should represent at least
one cycle across the image, $|b_{00}|/2^{k-1} \ge 1/x_\text{image}$---is
looser than \eqref{eq:k-max} in oversampled-detector geometries
typical of CT and is therefore not binding. The detector criterion is
the operative one.
\end{remark}

\subsection{Dual step scaling}\label{sec:sigma}

The ramp filter $\sqrt{|u|}$ places most of its energy at high
frequencies. Direct integration of $|f_i|^2$ shows that the dyadic
bands carry geometrically decreasing energy:
\begin{equation}\label{eq:band-energy}
  \int |f_i(u)|^2\,du \approx 2^{-i} \cdot \int |f_0(u)|^2\,du,
\end{equation}
so the LF tail contains a vanishingly small fraction of the total
spectral energy. Naively applying a single $\sigma$ to all bands
would mean the LF dual receives a comparatively tiny update per
iteration.

To compensate, we adopt a \emph{dyadic dual schedule}
\begin{equation}\label{eq:sigma-schedule}
  \sigma_i = 2^i \sigma_0,
\end{equation}
which exactly cancels the energy decay in \eqref{eq:band-energy}, so
that the contribution of band $i$ to the squared weighted operator
norm is approximately independent of $i$:
\begin{equation}\label{eq:sigma-cancels}
  \sigma_i \, \|R_i X\|^2 \approx 2^i \cdot 2^{-i} \cdot \sigma_0 \|R_0 X\|^2
    = \sigma_0 \|R_0 X\|^2.
\end{equation}
Consequently the weighted operator norm grows linearly in $k$ rather
than exponentially:
\begin{equation}\label{eq:norm-linear}
  \|\Sigma^{1/2} K_\text{data}\|^2
    = \sum_{i=0}^{k-1} \sigma_i \|R_i X\|^2
    \approx k \cdot \sigma_0 \|R_0 X\|^2.
\end{equation}
This is the key quantitative property of the dyadic split: it permits
$k$ independent dual variables at a cost of only an $O(k)$ factor in
the CP step constraint, which is much more favorable than the $O(2^k)$
factor one would incur with $\sigma_i = \sigma_0$ for all $i$ (because
of the unequal band energies).

\subsection{Saddle-point problem and convergence condition}

The $k$-channel saddle-point problem is
\begin{equation}\label{eq:cp-multi}
  \min_{f \ge 0}\ \alpha\, \mathrm{TV}_\text{D}(f) + \beta\, \|f\|_1
    \quad \text{s.t.}\quad
    \|R_i(X f - g)\|_2 \le \eps_i \sqrt{M},
    \quad i = 0, \ldots, k-1.
\end{equation}
Stacking the data, gradient, and $\ell_1$ channels gives the full
operator
\begin{equation}\label{eq:K-full}
  K = \begin{bmatrix}
    R_0 X \\ R_1 X \\ \vdots \\ R_{k-1} X \\
    \nabla_x \\ \nabla_y \\ I
  \end{bmatrix},
  \qquad
  \Sigma = \diag\big(\sigma_0 I,\, \sigma_1 I,\, \ldots,\, \sigma_{k-1} I,
    \sigma_\text{TV} I, \sigma_\text{TV} I, \sigma_{\ell_1} I\big),
\end{equation}
and the product-space CP convergence condition is
\begin{equation}\label{eq:cp-multi-cond}
  \boxed{\ \tau\, \|\Sigma^{1/2} K\|_2^2 < 1\ }
\end{equation}
which expands to
\begin{equation}\label{eq:cp-multi-explicit}
  \tau\, \left\|
  \begin{bmatrix}
    \sqrt{\sigma_0}\, R_0 X \\
    \sqrt{\sigma_1}\, R_1 X \\
    \vdots \\
    \sqrt{\sigma_{k-1}}\, R_{k-1} X \\
    \sqrt{\sigma_\text{TV}}\, \nabla_x \\
    \sqrt{\sigma_\text{TV}}\, \nabla_y \\
    \sqrt{\sigma_{\ell_1}}\, I
  \end{bmatrix}
  \right\|_2^2 < 1.
\end{equation}
This is the direct $k$-channel analogue of the two-channel condition
\eqref{eq:two-cond}. In practice we estimate $\|\Sigma^{1/2} K\|$ by
power iteration with $\sigma_0$ factored out, then choose
$\sigma_0 = c_\text{bal}/\|\Sigma^{1/2} K / \sqrt{\sigma_0}\|$ for a
target step-balance parameter $c_\text{bal}$ and derive $\tau$ from
$\tau \sigma_0 \|\Sigma^{1/2} K / \sqrt{\sigma_0}\|^2 = (1+\delta)^{-1}$.

\subsection{Algorithm}

The $k$-channel CP iteration alternates a primal step on $f$ with $k$
independent dual updates on the band variables $y_0, \ldots, y_{k-1}$,
the DTV duals $p_x, p_y$, and the $\ell_1$ dual $q$. The dual updates
on the band variables are projections onto closed $\ell_2$ balls of
radii $\eps_i \sqrt{M}$. We use the over-relaxation form with
extrapolation parameter $\rho \in [1, 2)$.

\begin{algorithm}[h]
\caption{Dyadic Multi-Channel CP for DTV-Regularized Limited-Angle CT}
\label{alg:multichannel-cp}
\begin{algorithmic}[1]
\State \textbf{Input:} sinogram $g$, projector $X$, dyadic filters $\{R_i\}_{i=0}^{k-1}$, DTV weights $\alpha_x, \alpha_y$, $\ell_1$ weight $\beta$, tolerances $\{\eps_i\}$, iterations $N$, relaxation $\rho$
\State \textbf{Compute step sizes:} estimate $\|\Sigma^{1/2} K\|$ by power iteration; set $\sigma_i = 2^i \sigma_0$ per \eqref{eq:sigma-schedule}, $\tau$ so that \eqref{eq:cp-multi-explicit} holds with slack $\delta$.
\State \textbf{Initialize:} $f^{(0)} = 0$, $\bar f^{(0)} = 0$, $y_i^{(0)} = 0$ for $i = 0,\ldots,k-1$, $p_x^{(0)} = p_y^{(0)} = q^{(0)} = 0$.
\For{$n = 0, \ldots, N-1$}
    \State \emph{(Primal step)}
    \State $f^{(n+1)} \gets \big[f^{(n)} - \tau\big(\sum_{i=0}^{k-1} R_i^* X^* y_i^{(n)}
        + \nabla_x^* p_x^{(n)} + \nabla_y^* p_y^{(n)} + q^{(n)}\big)\big]_+$
    \State $\bar f^{(n+1)} \gets 2 f^{(n+1)} - f^{(n)}$
    \For{$i = 0, \ldots, k-1$}
        \State \emph{(Per-band dual update: project onto $\ell_2$ ball of radius $\sigma_i \nu_\text{sino} \eps_i \sqrt{M}$)}
        \State $r_i \gets R_i X \bar f^{(n+1)} - R_i g$
        \State $\tilde y_i \gets y_i^{(n)} + \sigma_i r_i$
        \State $y_i^{(n+1)} \gets \tilde y_i \cdot \min\!\big(1,\, \sigma_i \nu_\text{sino} \eps_i \sqrt{M}\,/\,\|\tilde y_i\|\big)$
    \EndFor
    \State \emph{(DTV duals)}
    \State $\tilde p_x \gets p_x^{(n)} + \sigma_\text{TV} \nabla_x \bar f^{(n+1)};\quad p_x^{(n+1)} \gets \alpha_x \tilde p_x / \max(|\tilde p_x|, \alpha_x)$
    \State $\tilde p_y \gets p_y^{(n)} + \sigma_\text{TV} \nabla_y \bar f^{(n+1)};\quad p_y^{(n+1)} \gets \alpha_y \tilde p_y / \max(|\tilde p_y|, \alpha_y)$
    \State \emph{($\ell_1$ dual)}
    \State $\tilde q \gets q^{(n)} + \sigma_{\ell_1}\, \beta\, \bar f^{(n+1)};\quad q^{(n+1)} \gets \beta \tilde q / \max(|\tilde q|, \beta)$
    \State \emph{(Over-relaxation, $\rho \approx 1.75$)}
    \State $f^{(n+1)} \gets f^{(n)} - \rho(f^{(n)} - f^{(n+1)});\ \text{(same for } y_i, p_x, p_y, q\text{)}$
\EndFor
\State \Return $\bar f^{(N)}$
\end{algorithmic}
\end{algorithm}

\subsection{Practical notes}

\paragraph{Per-band tolerances.}
The vector $(\eps_0, \ldots, \eps_{k-1})$ encodes per-band data-fidelity
trust. Setting all $\eps_i$ equal recovers a matched-tolerance
multichannel solver. Setting $\eps_{k-1} < \eps_0$ tightens the LF
data fit relative to HF, which moves the fixed point in the direction
of better LF accuracy at the cost of HF data tolerance. This is the
correct lever for ``boosting low frequencies''---changing the $\eps_i$
modifies the solution, whereas changing the $\sigma_i$ alone only
modifies the optimization trajectory and is constrained by
\eqref{eq:cp-multi-cond}.

\paragraph{Per-iteration cost.}
Each iteration costs $k$ projections + $k$ backprojections for the
data channels, plus the DTV and $\ell_1$ operators. Per-iteration cost
is approximately $k\times$ the single-channel cost.

\paragraph{Choice of $\rho$.}
We use $\rho = 1.75$, consistent with \citet{sidky2012}. Lower values
($\rho = 1$) recover the unrelaxed CP form and give smoother but
slower convergence; higher values risk oscillation.

%=====================================================================
%\bibliography{references}
%\bibliographystyle{abbrvnat}
% Inline references until references.bib is populated:

\begin{thebibliography}{9}

\bibitem[Chambolle and Pock(2011)]{chambolle2011}
A.~Chambolle and T.~Pock.
\newblock A first-order primal-dual algorithm for convex problems with
  applications to imaging.
\newblock \emph{Journal of Mathematical Imaging and Vision}, 40(1):120--145,
  2011.

\bibitem[Pock and Chambolle(2011)]{pock2011}
T.~Pock and A.~Chambolle.
\newblock Diagonal preconditioning for first order primal-dual algorithms in
  convex optimization.
\newblock In \emph{2011 International Conference on Computer Vision},
  1762--1769. IEEE, 2011.

\bibitem[Sidky et al.(2012)]{sidky2012}
E.~Y. Sidky, J.~H. J{\o}rgensen, and X.~Pan.
\newblock Convex optimization problem prototyping for image reconstruction in
  computed tomography with the Chambolle-Pock algorithm.
\newblock \emph{Physics in Medicine and Biology}, 57(10):3065--3091, 2012.

\bibitem[Sidky et al.(2014)]{sidky2014}
E.~Y. Sidky, P.~Kao, and X.~Pan.
\newblock Two-channel reconstruction for sparse-view tomography.
\newblock \emph{Proceedings of the Third International Conference on Image
  Formation in X-ray Computed Tomography}, 2014.

\bibitem[Rigie and La Riviere(2015)]{rigie2015}
D.~S. Rigie and P.~J. La~Riviere.
\newblock Joint reconstruction of multi-channel, spectral CT data via
  constrained total nuclear variation minimization.
\newblock \emph{Physics in Medicine and Biology}, 60(5):1741--1762, 2015.

\bibitem[Kim et al.(2017)]{kim2017}
K.~Kim, T.~Possemier, and X.~Pan.
\newblock Directional total-variation regularization for limited-angle
  tomographic reconstruction.
\newblock In \emph{Proc. SPIE Medical Imaging}, 2017.

\end{thebibliography}

\end{document}